\section{实现方案}
\subsection{作品简介}
\ContentSlot{系统定位与总体功能}{说明系统的使用对象、预期工作流程、输入输出及拟交付形态。}
\subsection{技术方案}
\FigureSlot{系统总体技术方案}{预留总体架构图，表达感知、计算、通信、软件部署及应用之间的关系。}
\subsection{方案细节}
\KeyPoint{\StoryThreeTitle}{\StoryThreeQuestion}
{\small\renewcommand{\arraystretch}{\CompactTableStretch}
\noindent\begin{tabularx}{\linewidth}{@{}P{.25\linewidth}Y@{}}
\toprule
硬件配置项 & 方案配置 \\
\midrule
处理器与板卡 & \ChipModel{}；\BoardModel \\
FPGA & \FPGAModel \\
采集器件 & \SensorModel{}；\CameraModel \\
存储与接口 & \StoragePlan{}；\HardwareInterfaces \\
\bottomrule
\end{tabularx}}
\KeyPoint{\StoryFourTitle}{\StoryFourQuestion}
{\small\renewcommand{\arraystretch}{\CompactTableStretch}
\noindent\begin{tabularx}{\linewidth}{@{}P{.25\linewidth}Y@{}}
\toprule
软件配置项 & 方案配置 \\
\midrule
实时系统与Linux & \RealTimeSystem{}；\LinuxDistribution \\
工具链与板级支持 & \Toolchain{}；\BSPPlan \\
驱动与推理框架 & \DriverPlan{}；\InferenceFramework \\
混合部署 & \HybridDeploymentFunction \\
\bottomrule
\end{tabularx}}
\KeyPoint{\StoryFiveTitle}{\StoryFiveQuestion}
{\small\renewcommand{\arraystretch}{\CompactTableStretch}
\noindent\begin{tabularx}{\linewidth}{@{}P{.25\linewidth}Y@{}}
\toprule
功能配置项 & 方案配置 \\
\midrule
传感采集 & \AcquisitionFunction \\
通信与数据中心 & \CommunicationFunction{}；\MobileDataCenterFunction \\
视觉与模型推理 & \VisionFunction{}；\InferenceFunction{}；\AIModel \\
\bottomrule
\end{tabularx}}
\subsection{命题响应效果}
{\small\renewcommand{\arraystretch}{\CompactTableStretch}
\noindent\begin{tabularx}{\linewidth}{@{}P{.18\linewidth}YY@{}}
\toprule
指标 & 目标配置 & 验证方法 \\
\midrule
时延 & \TargetLatency & \LatencyVerificationMethod \\
帧率 & \TargetFrameRate & \FrameRateVerificationMethod \\
准确性 & \TargetAccuracy & \AccuracyVerificationMethod \\
功耗与资源 & \TargetPower{}；\TargetResourceUse & \PowerVerificationMethod{}；\ResourceUseVerificationMethod \\
\bottomrule
\end{tabularx}}
\ContentSlot{响应效果与验证依据}{按命题要求对应说明预期响应、评价方法和后续验证记录。}
\Needspace{7\baselineskip}
\subsection{命题企业对方案评价}
\ContentSlot{企业评价}{预留企业对接反馈、意见及相关依据的位置。}
